\documentclass[lettersize,journal]{IEEEtran}
\usepackage{amsmath,amsfonts,amssymb}
\usepackage{algorithmic}
\usepackage{algorithm}
\usepackage{array}
\usepackage[caption=false,font=normalsize,labelfont=sf,textfont=sf]{subfig}
\usepackage{textcomp}
\usepackage{stfloats}
\usepackage{url}
\usepackage{verbatim}
\usepackage{graphicx}
\usepackage{cite}
\usepackage{listings}
\usepackage{xcolor}
\usepackage[catalan]{babel}
\usepackage[T1]{fontenc}
\usepackage[utf8]{inputenc}
\usepackage{booktabs}
\usepackage{cuted}

% Configuracio de listings per a Java
\lstset{
    language=Java,
    basicstyle=\ttfamily\footnotesize,
    keywordstyle=\color{blue}\bfseries,
    commentstyle=\color{gray}\itshape,
    stringstyle=\color{red!70!black},
    numbers=left,
    numberstyle=\tiny\color{gray},
    numbersep=5pt,
    frame=single,
    breaklines=true,
    showstringspaces=false,
    tabsize=4,
    captionpos=b
}

\hyphenation{a-sin-to-tic a-sin-to-tics al-go-ris-me al-go-ris-mes con-tro-la-dor mul-ti-pli-ca-ti-va}

\begin{document}

\title{Anàlisi de Costos Computacionals Asimptòtics \\ Pràctica 1}

\author{Serafí~Nebot~Ginard}
% \thanks{Universitat de les Illes Balears (UIB), Grau en Enginyeria Informàtica, Anàlisi d'Algorismes, Curs 2025--2026.}}

\markboth{Pràctica 1, Curs 2025--2026}
{Serafí Nebot Ginard: Anàlisi de Costos Computacionals Asimptòtics}

\maketitle

\begin{strip}
\begin{abstract}
Aquesta memòria presenta el desenvolupament d'una aplicació Java 
que permet visualitzar i comparar el cost computacional asimptòtic de
quatre algorismes amb complexitats computacionals $O(n)$, $O(n \cdot \log n)$, $O(n^2)$
i $O(n^3)$. L'aplicació segueix estrictament el patró arquitectònic
Model-Vista-Controlador (MVC),
garantint una separació neta entre la lògica de negoci, la
interfície gràfica i les dades. Es presenten els fonaments
teòrics de l'anàlisi asimptòtic, el disseny del programari, la
identificació de la instrucció crítica de cada algorisme, el càlcul de la
constant multiplicativa i la predicció de temps d'execució per a mides
d'entrada arbitraries.
\end{abstract}

\end{strip}

%% ================================================================
\section{Introducció}
\IEEEPARstart{L}{'anàlisi} d'algorismes es una de les disciplines fonamentals
de la informàtica. Permet avaluar l'eficiència d'un algorisme en funció de
la mida de l'entrada, independentment del maquinari o del llenguatge de
programació emprat. El concepte central es el cost asimptòtic, que
descriu com creix el temps d'execució quan la mida de l'entrada $n$ tendeix
a infinit.

Quan es dissenya un algorisme, sovint existeixen múltiples estratègies per
resoldre el mateix problema. L'anàlisi asimptòtic proporciona una eina
objectiva per comparar-les: un algorisme amb cost $O(n \log n)$ sera
sempre mes eficient que un amb cost $O(n^2)$ per a entrades prou grans,
independentment de les constants implicades. Aquesta observació, que pot
semblar abstracta, te conseqüències practiques immediates: la diferència
entre un algorisme lineal i un cubic pot ser la diferencia entre una
execució de mil·lisegons i una de dies.

L'objectiu d'aquesta pràctica es doble. D'una banda, es proposa
implementar i mesurar experimentalment el temps d'execució de quatre
algorismes amb costos asimptòtics coneguts: lineal, quasilineal, quadratic
i cubic. D'altra banda, l'aplicació s'ha de construir seguint el patró
Model-Vista-Controlador (MVC) amb comunicació basada en el patró per
esdeveniments, un principi de disseny que desacobla completament les
responsabilitats de cada capa del programari.

El resultat es una aplicació d'escriptori Java Swing que permet a l'usuari
seleccionar rangs de mida d'entrada, executar algorismes de forma
individual o conjunta, visualitzar els resultats en una gràfica en temps
real (amb escala lineal o logarítmica), consultar les constants
multiplicatives mitjanes i fer previsions de temps per a mides d'entrada
no mesurades.

La memòria s'estructura de la següent manera: la Secció~II descriu
l'entorn de programació; la Secció~III exposa els fonaments teòrics; la
Secció~IV detalla el disseny i la implementació de la practica; la
Secció~V analitza el calcul del cost asimptòtic per la instrucció crítica;
la Secció~VI presenta els resultats; i la Secció~VII tanca amb les
conclusions.

%% ================================================================
% TODO: is this seciton necessary?
\section{Entorn de Programació}

L'aplicacio s'ha desenvolupat integrament en \textbf{Java}, un llenguatge
orientat a objectes que ofereix portabilitat, tipatge estatic i un ecosistema
madur de biblioteques. L'entorn concret es el seguent:

\begin{itemize}
    \item \textbf{Llenguatge:} Java 16 (compilacio amb \texttt{source} i
          \texttt{target} 16).
    \item \textbf{JDK:} OpenJDK 23 (Valhalla), compatible cap enrere amb
          Java~16.
    \item \textbf{IDE:} IntelliJ IDEA, que proporciona eines d'anàlisi de
          codi, refactoritzacio automatitzada i integracio amb Maven.
    \item \textbf{Sistema de construccio:} Apache Maven 3.x, amb un fitxer
          \texttt{pom.xml} que defineix el \texttt{groupId}
          \texttt{com.serafinebot.p1}, la codificacio UTF-8 i les versions
          del compilador.
    \item \textbf{Interficie grafica:} Java Swing (paquets \texttt{javax.swing}
          i \texttt{java.awt}), sense cap biblioteca externa de grafiques.
    \item \textbf{Control de versions:} Git.
\end{itemize}

La decisio d'utilitzar exclusivament Java Swing ---sense biblioteques
de tercers com JFreeChart--- respon a un requisit de la practica i obliga
a implementar manualment el renderitzat de la grafica amb l'API
\texttt{Graphics2D}, incloent-hi eixos, graella, series de dades i
llegenda.

El fitxer \texttt{pom.xml} segueix una configuracio minima, ja que el
projecte no te cap dependencia externa:

\begin{lstlisting}[caption={Configuracio Maven del projecte.},language=XML]
<properties>
  <maven.compiler.source>16</maven.compiler.source>
  <maven.compiler.target>16</maven.compiler.target>
  <project.build.sourceEncoding>
    UTF-8
  </project.build.sourceEncoding>
</properties>
\end{lstlisting}

Aquesta absencia de dependencies externes es un aspecte rellevant del
projecte: tota la funcionalitat ---des del renderitzat de grafiques fins a
la concurrencia--- s'implementa amb les eines estandard del JDK.

%% ================================================================
\section{Fonaments Teòrics}

\subsection{Anàlisi Asimptòtic de Costos}

L'anàlisi asimptòtic permet classificar algorismes segons el seu
comportament a llarg termini, es a dir, com creix el seu temps d'execució
en funció de la mida de l'entrada~\cite{ref_cormen}. Es defineixen les
notacions següents:

\begin{itemize}
    \item \textbf{$O(f(n))$} (cota superior asimptòtica): existeixen
          constants $c > 0$ i $n_0$ tals que $T(n) \leq c \cdot f(n)$ per
          a tot $n \geq n_0$.
    \item \textbf{$\Omega(f(n))$} (cota inferior): existeixen constants
          $c > 0$ i $n_0$ tals que $T(n) \geq c \cdot f(n)$ per a tot
          $n \geq n_0$.
    \item \textbf{$\Theta(f(n))$} (cota ajustada): $T(n)$ es alhora
          $O(f(n))$ i $\Omega(f(n))$.
\end{itemize}

El temps d'execució real d'un algorisme es pot modelar com:
\begin{equation}
\label{eq:temps}
T(n) = c \cdot f(n)
\end{equation}
on $f(n)$ es la funció de cost teòric i $c$ es la \textit{constant
multiplicativa}, que depèn del maquinari, el compilador i altres factors
constants. La notació asimptòtica extreu aquesta constant per a centrar-se només
en l'ordre de creixement.

\subsection{Ordres de Creixement}

Les quatre funcions de cost que s'analitzen en aquesta pràctica presenten
comportaments de creixement dràsticament diferents. La Taula~\ref{tab:growth}
il·lustra aquesta diferència amb valors numèrics concrets.

\begin{table}[!t]
      \caption{Comparació d'ordres de creixement per a diferents valors de $n$.\label{tab:growth}}
      \centering
      \begin{tabular}{r|rrrr}
            \toprule
            $n$ & $n$ & $n \log_2 n$ & $n^2$ & $n^3$ \\
            \midrule
            10     & 10       & 33         & 100         & 1\,000 \\
            100    & 100      & 664        & 10\,000     & $10^6$ \\
            1\,000  & 1\,000  & 9\,966     & $10^6$      & $10^9$ \\
            10\,000 & 10\,000 & 132\,877   & $10^8$      & $10^{12}$ \\
            \bottomrule
      \end{tabular}
\end{table}

Com es pot observar, per a $n = 10\,000$, un algorisme cubic realitza
$10^{12}$ operacions basiques, mentre que un de lineal nomes en fa
$10^4$: una diferencia de vuit ordres de magnitud. Si cada operacio
trigues 1~nanosegon, l'algorisme lineal acabaria en 10~microsegons,
mentre que el cubic necessitaria aproximadament 17~minuts.

\subsection{Instruccio Critica}

La \textit{instruccio critica} es aquella operacio dins l'algorisme que
s'executa el major nombre de vegades i que, per tant, domina el cost
total~\cite{ref_knuth}. Identificar-la es clau per determinar la
complexitat asimptòtica. En algorismes iteratius, la instruccio critica
es habitualment l'operacio dins el bucle mes intern. El nombre total
d'execucions d'aquesta instruccio depen del nombre d'iteracions de tots
els bucles que l'envolten.

Per exemple, en un algorisme amb dos bucles imbricats, cadascun iterant
$n$ vegades, la instruccio critica s'executa $n \times n = n^2$ vegades,
i per tant el cost asimptòtic es $\Theta(n^2)$.

\subsection{Patro Model-Vista-Controlador (MVC)}

El patro MVC~\cite{ref_mvc} divideix una aplicacio en tres components
amb responsabilitats ben definides:

\begin{enumerate}
    \item \textbf{Model:} encapsula les dades i la logica de negoci.
          Notifica als observadors quan el seu estat canvia, pero no coneix
          la interficie grafica. En el nostre cas, emmagatzema les mesures
          d'execucio, calcula constants multiplicatives i fa prediccions.
    \item \textbf{Vista:} es responsable de la presentacio visual.
          Renderitza les dades del model i captura les accions de l'usuari,
          pero no coneix la logica de negoci. En el nostre cas, dibuixa la
          grafica, mostra la taula de resultats i ofereix els controls.
    \item \textbf{Controlador:} actua com a intermediari. Rep events
          semantics de la vista, invoca operacions sobre el model i
          actualitza la vista amb resultats. En el nostre cas, gestiona
          els fils d'execucio dels algorismes.
\end{enumerate}

Un principi fonamental del MVC ben implementat es que la \textbf{Vista mai
no ha d'exposar els seus components visuals} (botons, camps de text, etc.)
al Controlador ni a cap altra capa. Tota comunicacio ha de ser a traves
d'events semantics i metodes que operen amb dades. Aquesta restriccio es
crucial perque, si el controlador accedis directament als botons o camps
de text de la vista, qualsevol canvi a la interficie grafica obligaria a
modificar tambe el controlador, trencant el principi de separacio de
responsabilitats.

En implementacions ingenues del MVC, es comu veure patrons com
\texttt{view.getStartButton().addActionListener(...)}, on el controlador
obte directament un \texttt{JButton} de la vista. Aixo viola el patro
perque crea un acoblament directe entre el controlador i la representacio
visual concreta. En la nostra implementació, hem evitat completament
aquest anti-patro.

\subsection{Patro per Esdeveniments (Observer)}

El patro Observer~\cite{ref_observer} estableix una relacio
un-a-molts entre objectes, de manera que quan un objecte (el
\textit{subjecte}) canvia d'estat, tots els seus observadors son
notificats automaticament. En el context d'aquesta practica, s'utilitza
en dues direccions:

\begin{enumerate}
    \item \textbf{Model $\rightarrow$ Vista:} El model defineix una
          interficie \texttt{ModelListener} amb el metode
          \texttt{modelChanged()}. Quan s'afegeix una mesura o s'esborren
          dades, el model invoca aquest metode, i la vista actualitza la
          taula i la grafica.
    \item \textbf{Vista $\rightarrow$ Controlador:} La vista defineix
          una interficie \texttt{ViewListener} amb set metodes semantics
          (\texttt{onStartAlgorithm}, \texttt{onStopAlgorithm},
          \texttt{onStartAll}, \texttt{onStopAll}, \texttt{onPredict},
          \texttt{onClear}, \texttt{onLogScaleChanged}). El controlador
          implementa aquesta interficie i rep events ja processats (amb
          valors parsejats), mai components Swing.
\end{enumerate}

Aquesta doble aplicacio del patro Observer garanteix un desacoblament
complet: el model no sap qui l'observa, i el controlador no sap com es
la interficie grafica. Es podria substituir tota la vista Swing per una
interficie web o de terminal sense modificar ni una sola linia del model
o del controlador.

\subsection{Concurrencia i el Fil EDT}

En Java Swing, totes les operacions sobre components visuals s'han de
realitzar des del fil de despatx d'events (\textit{Event Dispatch Thread},
EDT)~\cite{ref_java}. Si un algorisme lent s'executes directament al fil
EDT, la interficie es bloquejaria fins que acabes, impedint qualsevol
interaccio amb l'usuari.

La solucio estandard es \texttt{SwingWorker}, una classe que permet
executar tasques llargues en un fil de fons i publicar resultats parcials
al fil EDT de forma segura. El cicle de vida d'un \texttt{SwingWorker} es:

\begin{enumerate}
    \item \texttt{doInBackground()}: s'executa en un fil de fons. Aqui
          s'executa l'algorisme i es crida \texttt{publish()} per enviar
          resultats intermedis.
    \item \texttt{process(List<V>)}: s'invoca al fil EDT amb els resultats
          publicats. Aqui s'afegeixen les mesures al model.
    \item \texttt{done()}: s'invoca al fil EDT quan la tasca acaba o es
          cancel-la. Aqui s'actualitza l'estat visual del boto.
\end{enumerate}

La cancel-lacio es cooperativa: \texttt{SwingWorker.cancel(true)} envia
una interrupcio al fil de fons, que ha de comprovar periodicament
\texttt{Thread.currentThread().isInterrupted()} per aturar-se netament.

%% ================================================================
\section{Disseny i Implementació de la Practica}

\subsection{Estructura del Projecte}

El projecte Maven s'organitza en tres paquets dins
\texttt{com.serafinebot.p1}:

\begin{verbatim}
src/main/java/com/serafinebot/p1/
  Main.java
  model/
    AlgorithmType.java
    Measurement.java
    Model.java
  view/
    ViewListener.java
    ControlPanel.java
    GraphPanel.java
    ResultsPanel.java
    MainView.java
  controller/
    Controller.java
\end{verbatim}

Cada paquet correspon exactament a una capa del patro MVC, amb el punt
d'entrada \texttt{Main.java} a l'arrel. L'ordre d'instanciacio segueix
la sequencia classica:

\begin{lstlisting}[caption={Punt d'entrada de l'aplicacio (\texttt{Main.java}).}]
Model model = new Model();
MainView view = new MainView(model);
new Controller(model, view);
view.setVisible(true);
\end{lstlisting}

El model es crea primer perque no depen de res; la vista rep el model per
poder registrar-se com a observadora dels seus canvis; i el controlador
rep ambdos i es registra com a \texttt{ViewListener} de la vista. Tota
la inicialitzacio es realitza dins \texttt{SwingUtilities.invokeLater()}
per garantir que es fa al fil EDT.

\subsection{Capa del Model}

La capa del model consta de tres classes, cadascuna amb una responsabilitat
clara i ben delimitada.

\subsubsection{AlgorithmType (Enum)}

Defineix els quatre algorismes com a constants d'un \texttt{enum}. Cada
constant implementa dos metodes abstractes:
\begin{itemize}
    \item \texttt{execute(long n)}: executa l'algorisme amb mida $n$ i
          retorna un resultat numericament dependent de l'entrada per
          evitar que el compilador JIT l'optimitzi o l'elimini.
    \item \texttt{theoreticalCost(long n)}: retorna el valor numeric de
          $f(n)$.
\end{itemize}

A mes, cada algorisme te associat un nom visible, una notacio matematica
i un color per a la grafica. L'us d'un \texttt{enum} amb metodes
abstractes es una aplicacio del \textit{Strategy pattern}: cada constant
de l'enumeracio encapsula el seu propi comportament.

Per exemple, l'algorisme quadratic es defineix com:

\begin{lstlisting}[caption={Definicio de l'algorisme $O(n^2)$.}]
QUADRATIC("O(n^2)", "Quadratic",
          new Color(220, 120, 0)) {
    public double theoreticalCost(long n) {
        return (double) n * n;
    }
    public long execute(long n)
            throws InterruptedException {
        long sum = 0;
        for (long i = 0; i < n; i++) {
            if (Thread.currentThread()
                      .isInterrupted())
                throw new InterruptedException();
            for (long j = 0; j < n; j++) {
                sum += i * j; // instruccio critica
            }
        }
        return sum;
    }
};
\end{lstlisting}

Cal destacar dos aspectes: primer, la comprovacio d'\texttt{isInterrupted()}
al bucle extern (no al mes intern, per no penalitzar el rendiment amb
comprovacions massa frequents); segon, el llancament d'\texttt{InterruptedException}
que permet la cancel-lacio cooperativa des del controlador.

\subsubsection{Measurement}

Classe immutable que emmagatzema el resultat d'una mesura: el tipus
d'algorisme, la mida $n$, el temps en nanosegons i la constant
multiplicativa, que es pre-calcula al constructor:

\begin{equation}
\label{eq:constant}
c = \frac{T(n)}{f(n)}
\end{equation}

on $T(n)$ es el temps mesurat en nanosegons i $f(n)$ es el cost teòric
retornat per \texttt{theoreticalCost(n)}. El codi del constructor es:

\begin{lstlisting}[caption={Constructor de \texttt{Measurement}.}]
public Measurement(AlgorithmType type,
                   long n, long timeNanos) {
    this.type = type;
    this.n = n;
    this.timeNanos = timeNanos;
    double cost = type.theoreticalCost(n);
    this.constant = cost > 0
        ? timeNanos / cost : 0;
}
\end{lstlisting}

Pre-calcular $c$ a cada mesura individual permet despres obtenir la
mitjana de forma directa, sense necessitat de recalcular divisions
repetidament.

\subsubsection{Model}

Classe central que manté un \texttt{EnumMap} de llistes de mesures. Els
metodes principals son:
\begin{itemize}
    \item \texttt{addMeasurement(Measurement m)}: afegeix una mesura,
          ordena la llista per $n$ creixent (per facilitar el dibuix de
          la grafica) i notifica els observadors.
    \item \texttt{getAverageConstant(AlgorithmType t)}: retorna la mitjana
          aritmetica de les constants individuals, calculada amb l'API
          de \textit{streams} de Java.
    \item \texttt{predict(AlgorithmType t, long n)}: calcula el temps
          predit com $\hat{T}(n) = \bar{c} \cdot f(n)$, on $\bar{c}$ es
          la constant mitjana.
    \item \texttt{clearAll()}: esborra totes les mesures i notifica.
\end{itemize}

Tots els metodes que modifiquen o llegeixen l'estat son
\texttt{synchronized}, ja que els fils de fons dels \texttt{SwingWorker}
poden accedir al model concurrentment. El model implementa el patro
Observer mitjancant la interficie \texttt{ModelListener}:

\begin{lstlisting}[caption={Patro Observer al model.}]
public interface ModelListener {
    void modelChanged();
}

private void fireModelChanged() {
    for (ModelListener l : listeners) {
        l.modelChanged();
    }
}
\end{lstlisting}

\subsection{Capa de la Vista}

La vista es composa de quatre classes i una interficie. \textbf{Cap
d'elles exposa components Swing a l'exterior.}

\subsubsection{ViewListener (Interficie)}

Defineix el contracte semantic entre la vista i el controlador:

\begin{lstlisting}[caption={Interficie \texttt{ViewListener}.}]
public interface ViewListener {
    void onStartAlgorithm(AlgorithmType type,
                          long[] nValues);
    void onStopAlgorithm(AlgorithmType type);
    void onStartAll(long[] nValues);
    void onStopAll();
    void onPredict(long predictN);
    void onClear();
    void onLogScaleChanged(boolean logScale);
}
\end{lstlisting}

Tots els parametres son dades primitives o enums. No hi ha cap referencia
a \texttt{JButton}, \texttt{JTextField} ni cap altre component visual.
Aquesta es la clau de la separacio MVC.

\subsubsection{ControlPanel}

Panell que conte els camps d'entrada, els botons d'execucio i el camp de
previsio. Les responsabilitats son exclusivament de la vista:

\begin{enumerate}
    \item Capturar les accions de l'usuari (clics de botons, canvi de
          checkbox d'escala logaritmica).
    \item Validar l'entrada internament: parsejar els camps de text a
          valors \texttt{long}/\texttt{int} i mostrar dialogs d'error
          (\texttt{JOptionPane}) si els valors son invalids.
    \item Generar l'array de valors de $n$ equidistants dins el rang
          $[\texttt{nInicial}, \texttt{nFinal}]$ amb la formula:
\end{enumerate}

\begin{equation}
n_i = n_{\text{inici}} + \frac{(n_{\text{final}} - n_{\text{inici}}) \cdot i}{p - 1}, \quad i = 0, \ldots, p-1
\end{equation}

on $p$ es el nombre de passos. Nomes quan la validacio passa
satisfactoriament s'emet l'event semantic a traves de \texttt{ViewListener}.

Quan un algorisme esta en execucio, el seu boto canvia d'icona
($\blacktriangleright \rightarrow \blacksquare$) i de color (del color
de l'algorisme a vermell), permetent aturar-lo individualment amb un
segon clic.

\subsubsection{GraphPanel}

Panell personalitzat que sobreescriu \texttt{paintComponent(Graphics g)}
per dibuixar la grafica amb l'API \texttt{Graphics2D}. Les
caracteristiques principals son:

\begin{itemize}
    \item \textbf{Eixos i graella:} eix X per a $n$, eix Y per al temps
          en mil-lisegons, amb marques numeriques i linies de graella
          de color gris clar.
    \item \textbf{Escala dual:} suporta escala lineal i logaritmica a
          l'eix Y, commutable en temps real. En escala logaritmica, la
          posicio vertical es calcula com:
\end{itemize}

\begin{equation}
y = h - h \cdot \frac{\log_{10}(t) - \log_{\min}}{\log_{\max} - \log_{\min}}
\end{equation}

\begin{itemize}
    \item \textbf{Series de dades:} dibuixa punts (cercles plens de
          8$\times$8 pixels) i linies de connexio per a cada algorisme
          amb el seu color assignat i un tracat de 2.5~punts.
    \item \textbf{Llegenda:} rectangle semitransparent amb fons blanc i
          opacitat 220/255 a la cantonada superior esquerra.
    \item \textbf{Formatejat dinamic:} les etiquetes s'adapten
          automaticament (ns, $\mu$s, ms, s, min per al temps; K, M per
          a $n$).
    \item \textbf{Antialiasing:} activat tant per a linies com per a
          text per obtenir un renderitzat suau.
\end{itemize}

L'escala logaritmica es especialment util quan es comparen algorismes
amb ordres de magnitud molt diferents (per exemple, $O(n)$ vs. $O(n^3)$),
ja que en escala logaritmica les funcions potencials $n^k$ es visualitzen
com rectes amb pendent $k$.

\subsubsection{ResultsPanel}

Panell dividit verticalment en dues parts mitjancant un \texttt{JSplitPane}:
\begin{itemize}
    \item Una \texttt{JTable} no editable que mostra totes les mesures
          amb quatre columnes: algorisme, $n$, temps en ms i constant
          $c$ en notacio cientifica.
    \item Una \texttt{JTextArea} de nomes lectura amb font monoespaiada
          que mostra les constants multiplicatives mitjanes i les
          previsions sol-licitades per l'usuari, amb formatejat adaptatiu
          (mil-lisegons, segons, minuts, hores o dies).
\end{itemize}

\subsubsection{MainView}

\texttt{JFrame} principal que compon els tres sub-panells en un
\texttt{BorderLayout} (controls al \texttt{NORTH}, grafica al
\texttt{CENTER}, resultats a l'\texttt{EAST}). Es registra com a
observadora del model amb una lambda:

\begin{lstlisting}[caption={Registre com a observadora del model.}]
model.addListener(() -> {
    graphPanel.repaint();
    resultsPanel.updateResults();
});
\end{lstlisting}

L'API publica de \texttt{MainView} consta exclusivament de metodes
que operen amb dades:

\begin{lstlisting}[caption={API publica de \texttt{MainView}.}]
void addViewListener(ViewListener l);
void setAlgorithmRunning(
    AlgorithmType t, boolean b);
void setLogScale(boolean logScale);
void showPredictions(long n,
    Map<AlgorithmType, Double> preds);
\end{lstlisting}

Cap metode retorna ni rep un component Swing. Aquesta restriccio es
fonamental per a la separacio MVC i ha estat un dels punts de disseny
mes importants del projecte.

\subsection{Capa del Controlador}

La classe \texttt{Controller} implementa \texttt{ViewListener} i gestiona
el cicle de vida dels algorismes. El codi del metode principal d'execucio
mostra el patro de \texttt{SwingWorker}:

\begin{lstlisting}[caption={Execucio d'un algorisme al controlador.}]
private void startExecution(
        AlgorithmType type, long[] nValues) {
  SwingWorker<Void, Measurement> worker =
    new SwingWorker<>() {
      protected Void doInBackground() {
        for (long n : nValues) {
          if (isCancelled()) break;
          // Escalfament JIT
          type.execute(Math.min(n, 500));
          // Mesura real
          long start = System.nanoTime();
          type.execute(n);
          long elapsed = System.nanoTime()-start;
          publish(new Measurement(
              type, n, elapsed));
        }
        return null;
      }
      protected void process(
              List<Measurement> chunks) {
        for (Measurement m : chunks)
          model.addMeasurement(m);
      }
      protected void done() {
        workers.remove(type);
        view.setAlgorithmRunning(type, false);
      }
    };
  workers.put(type, worker);
  view.setAlgorithmRunning(type, true);
  worker.execute();
}
\end{lstlisting}

Cal destacar l'\textit{escalfament} (warm-up) que es fa abans de cada
mesura: s'executa l'algorisme amb una mida petita ($\min(n, 500)$) per
forcar la compilacio JIT de la JVM. Aixo evita que la primera mesura
estigui distorsionada pel cost de compilacio.

El controlador mante un \texttt{EnumMap} de \texttt{SwingWorker} per
poder cancel-lar cada algorisme individualment. \textbf{Mai no accedeix a
cap component visual}: la seva unica referencia a la vista es a traves
dels metodes de dades de \texttt{MainView}.

%% ================================================================
\section{Calcul del Cost Asimptòtic\\ per la Instruccio Critica}

La instruccio critica de cada algorisme es l'operacio dins el bucle mes
intern que s'executa el major nombre de vegades. El nombre total
d'execucions d'aquesta instruccio determina $f(n)$.

\subsection{Algorisme Lineal --- $O(n)$}

\begin{lstlisting}[caption={Algorisme lineal.}]
for (long i = 0; i < n; i++) {
    sum += i * i;  // instruccio critica
}
\end{lstlisting}

La instruccio critica \texttt{sum += i * i} s'executa exactament $n$
vegades. No hi ha bucles imbricats, per tant:
\begin{equation}
f(n) = n \implies T(n) \in \Theta(n)
\end{equation}

\subsection{Algorisme Quasilineal --- $O(n \cdot \log n)$}

\begin{lstlisting}[caption={Algorisme $O(n \cdot \log n)$.}]
int logN = (int)(Math.log(n) / Math.log(2));
for (long i = 0; i < n; i++) {
    for (int j = 0; j < logN; j++) {
        sum += i ^ j;  // instruccio critica
    }
}
\end{lstlisting}

El bucle extern itera $n$ vegades i l'intern $\lfloor \log_2 n \rfloor$
vegades per cada iteracio externa. La instruccio critica
\texttt{sum += i \^{} j} s'executa:
\begin{equation}
f(n) = n \cdot \lfloor \log_2 n \rfloor \implies T(n) \in \Theta(n \log n)
\end{equation}

\subsection{Algorisme Quadratic --- $O(n^2)$}

\begin{lstlisting}[caption={Algorisme quadratic.}]
for (long i = 0; i < n; i++) {
    for (long j = 0; j < n; j++) {
        sum += i * j;  // instruccio critica
    }
}
\end{lstlisting}

Dos bucles imbricats, cadascun de $n$ iteracions. La instruccio critica
\texttt{sum += i * j} s'executa:
\begin{equation}
f(n) = n \times n = n^2 \implies T(n) \in \Theta(n^2)
\end{equation}

\subsection{Algorisme Cubic --- $O(n^3)$}

\begin{lstlisting}[caption={Algorisme cubic.}]
for (long i = 0; i < n; i++) {
    for (long j = 0; j < n; j++) {
        for (long k = 0; k < n; k++) {
            sum += (i*j) + k; // inst. critica
        }
    }
}
\end{lstlisting}

Tres bucles imbricats, cadascun de $n$ iteracions. La instruccio critica
\texttt{sum += (i * j) + k} s'executa:
\begin{equation}
f(n) = n^3 \implies T(n) \in \Theta(n^3)
\end{equation}

\subsection{Resum Comparatiu}

La Taula~\ref{tab:critical} resumeix la instruccio critica, el nombre
d'execucions i la complexitat de cada algorisme.

\begin{table}[!t]
\caption{Instruccio critica i complexitat de cada algorisme.\label{tab:critical}}
\centering
\begin{tabular}{lllc}
\toprule
Algorisme & Instruccio critica & Execucions & $\Theta$ \\
\midrule
Lineal     & \texttt{sum += i*i}   & $n$                           & $n$ \\
N Log N    & \texttt{sum += i\^{}j}  & $n \lfloor\log_2 n\rfloor$   & $n\log n$ \\
Quadratic  & \texttt{sum += i*j}   & $n^2$                         & $n^2$ \\
Cubic      & \texttt{sum += (i*j)+k} & $n^3$                       & $n^3$ \\
\bottomrule
\end{tabular}
\end{table}

\subsection{Constant Multiplicativa i Prediccio}

Per a cada mesura individual, la constant multiplicativa es calcula
segons l'equacio~(\ref{eq:constant}). La constant mitjana $\bar{c}$
s'obte com la mitjana aritmetica de totes les constants individuals
d'un algorisme:

\begin{equation}
\bar{c} = \frac{1}{m} \sum_{k=1}^{m} c_k = \frac{1}{m} \sum_{k=1}^{m}
\frac{T(n_k)}{f(n_k)}
\end{equation}

on $m$ es el nombre de mesures realitzades. Aquesta constant reflexa el
cost temporal per unitat d'operacio basica a la maquina concreta.

Per fer una prediccio del temps d'execucio per a un valor de $n$ no
mesurat, s'utilitza la formula:
\begin{equation}
\hat{T}(n) = \bar{c} \cdot f(n)
\end{equation}

Per exemple, si tenim $\bar{c} = 2.5$ ns per a l'algorisme quadratic
i volem predir $n = 100\,000$:
\begin{equation}
\hat{T}(100\,000) = 2.5 \times (100\,000)^2 = 2.5 \times 10^{10} \text{ ns}
= 25 \text{ s}
\end{equation}

Aquesta capacitat de prediccio es un dels avantatges clau de l'anàlisi
asimptòtic: amb unes poques mesures per a valors petits de $n$, es pot
estimar el temps d'execucio per a valors molt mes grans sense necessitat
d'executar l'algorisme.

%% ================================================================
\section{Resultats i Discussio}

L'aplicacio implementada permet verificar experimentalment que els temps
d'execucio mesurats segueixen les funcions de cost teòric. A la grafica,
s'observa clarament com:

\begin{itemize}
    \item La serie lineal $O(n)$ creix de forma gairebe imperceptible en
          comparacio amb les altres series.
    \item La serie $O(n \cdot \log n)$ creix lleugerament mes que la
          lineal, pero de forma moderada i controlada.
    \item La serie quadratica $O(n^2)$ creix de forma parabolica,
          essent significativament mes costosa per a valors grans de $n$.
    \item La serie cubica $O(n^3)$ domina clarament totes les altres,
          fent-se impracticable per a valors de $n$ moderadament grans.
\end{itemize}

L'us de l'escala logaritmica a l'eix Y permet comparar totes quatre series
simultaniament sense que les mes rapides quedin comprimides contra l'eix X.
En escala logaritmica, les funcions potencials $n^k$ es visualitzen com a
rectes amb pendent proporcional a $k$, facilitant la verificacio visual
de la complexitat.

Les constants multiplicatives calculades es mantenen raonablement estables
per a un algorisme donat, la qual cosa confirma que el model asimptòtic
$T(n) = c \cdot f(n)$ es una bona aproximacio. Petites variacions
s'atribueixen a:

\begin{itemize}
    \item Efectes de la jerarquia de memoria (cache L1/L2/L3 del
          processador): per a valors grans de $n$, les dades poden
          excedir la cache, augmentant la latencia.
    \item La compilacio JIT de la JVM: el compilador optimitza el codi
          progressivament durant l'execucio, la qual cosa pot fer que
          les primeres mesures siguin mes lentes. L'escalfament previ
          mitiga parcialment aquest efecte.
    \item La planificacio del sistema operatiu: altres processos poden
          competir pel temps de CPU.
\end{itemize}

La funcionalitat de prediccio ha mostrat resultats coherents: les
prediccions per a valors de $n$ dins el rang mesurat concorden amb les
mesures reals, i les prediccions per a valors mes grans segueixen
l'extrapolacio esperada.

Respecte al disseny del programari, l'estricta separacio MVC amb
comunicacio per esdeveniments ha permes:

\begin{itemize}
    \item Desenvolupar i provar la logica del model independentment de la
          vista.
    \item Modificar la interficie grafica (per exemple, afegir l'escala
          logaritmica o millorar el formatejat de la grafica) sense alterar
          el controlador ni el model.
    \item Garantir que l'aplicacio roman responsiva durant l'execucio
          d'algorismes lents, gracies a l'us de \texttt{SwingWorker} que
          executa els algorismes en fils de fons separats del fil EDT.
    \item Cancel-lar algorismes individualment sense afectar els altres
          que s'estan executant concurrentment.
\end{itemize}

%% ================================================================
\section{Conclusions}

Aquesta practica ha permes aplicar de forma conjunta els conceptes
d'anàlisi asimptòtic i de disseny de programari amb el patro MVC.

Des del punt de vista teòric, s'ha verificat experimentalment que els
temps d'execucio dels quatre algorismes implementats corresponen a les
seves complexitats teoriques. La identificacio de la instruccio critica
de cada algorisme ---l'operacio del bucle mes intern--- i el calcul de la
constant multiplicativa han permes fer prediccions fiables del temps
d'execucio per a mides d'entrada arbitraries. L'experiment confirma que
el model $T(n) = c \cdot f(n)$ es una aproximacio valida quan $n$ es prou
gran i les condicions d'execucio son estables.

Des del punt de vista del disseny de programari, s'ha aconseguit una
implementació estricta del patro MVC on:

\begin{enumerate}
    \item La vista \textbf{no exposa cap component visual} al controlador;
          tota interaccio es a traves d'events semantics.
    \item La interficie \texttt{ViewListener} defineix set events amb
          parametres de dades (enums, arrays de \texttt{long}, booleans),
          mai objectes Swing.
    \item El model notifica canvis a traves de \texttt{ModelListener}
          sense coneixer qui l'observa.
    \item El controlador opera exclusivament amb dades, delegant la
          presentacio visual a la vista i la logica de negoci al model.
    \item La concurrencia es gestiona amb \texttt{SwingWorker}, garantint
          la responsivitat de la interficie i la cancel-lacio segura.
\end{enumerate}

El patro per esdeveniments ha demostrat ser essencial per mantenir el
desacoblament entre capes, permetent que cada component evolucioni de
forma independent. Aquesta separacio no es una mera formalitat academica:
facilita el manteniment, la testabilitat i l'extensibilitat del
programari. Per exemple, es podria substituir la vista Swing per una
interficie JavaFX o web sense modificar ni una linia del model o del
controlador.

Com a treball futur, es podrien afegir mes algorismes (per exemple,
$O(2^n)$ o $O(n!)$), implementar l'exportacio de dades a CSV, permetre
a l'usuari definir algorismes personalitzats, o afegir mes estrategies
de mesura com la repeticio multiple amb mitjana per reduir la variancia
experimental.

%% ================================================================

\begin{thebibliography}{9}
\bibliographystyle{IEEEtran}

\bibitem{ref_mvc}
E. Gamma, R. Helm, R. Johnson, and J. Vlissides, \textit{Design Patterns:
Elements of Reusable Object-Oriented Software}. Reading, MA, USA:
Addison-Wesley, 1994.

\bibitem{ref_observer}
F. Buschmann, R. Meunier, H. Rohnert, P. Sommerlad, and M. Stal,
\textit{Pattern-Oriented Software Architecture: A System of Patterns}.
Chichester, U.K.: Wiley, 1996.

\bibitem{ref_cormen}
T. H. Cormen, C. E. Leiserson, R. L. Rivest, and C. Stein,
\textit{Introduction to Algorithms}, 4th~ed. Cambridge, MA, USA:
MIT Press, 2022.

\bibitem{ref_sedgewick}
R. Sedgewick and K. Wayne, \textit{Algorithms}, 4th~ed. Upper Saddle
River, NJ, USA: Addison-Wesley, 2011.

\bibitem{ref_swing}
J. Elliott, R. Eckstein, M. Loy, D. Wood, and B. Cole, \textit{Java
Swing}, 2nd~ed. Sebastopol, CA, USA: O'Reilly Media, 2002.

\bibitem{ref_java}
Oracle Corporation, ``The Java Tutorials --- Concurrency,''
[Online]. Available:
\url{https://docs.oracle.com/javase/tutorial/essential/concurrency/}

\bibitem{ref_knuth}
D. E. Knuth, \textit{The Art of Computer Programming, Vol. 1:
Fundamental Algorithms}, 3rd~ed. Reading, MA, USA: Addison-Wesley, 1997.

\end{thebibliography}

\end{document}
